\scanpage{14}
\deleted{à un \unsure{gpe|groupe} fini.} Or un élément de
$J(k(\eta))$ \deleted{\unsure{provi|qui}} provient \deleted{\unsure{peut} \unsure{être} \unsure{de la} \unsure{une}
obstruction de \unsure{t}} d’une classe de
\unsure{diviseur|diviseurs} de degré $0$ sur $Y_\eta$

[N.B. le groupe de Brauer de $k(\eta)=k(\unsure{s|t})$
est nul, par Tsen] et on gagne

\[
 \boxed{\rho(\widetilde X)-\unsure{2|1}=}\boxed{\rho(X)+z-2=\rho(J')}.
\]

\[
\left\{\begin{aligned}
&\rho(X)\quad\text{nb de Picard de }X\\
&z=\text{nb de sections singulières du pinceau}\\
&\rho(J')\quad\deleted{\unsure{\mathrm{nb}|=}}\ \text{rang du groupe des pts rat. sur }k(\eta)\\
&\hspace{3em}\text{de }J'=J(X_\eta)/J(X_\eta)^{\rm fixe}
\end{aligned}\right.
\]

Condition de validité de la formule :
toutes les sections du pinceau sont \deleted{intègres.}
\uline{irréductibles} (intègre n’est pas absolument \unsure{néc|néc.},
i.e. on admet des sections multiples \ldots).
% archive sheet number: 13
